% GENERATED FILE. Edit canonical lessons, then run npm run generate:books.
% canonical-source-hash: fnv1a64:47b32990ae14b382
% canonical-lessons: JA-C236-eiga, JA-C236-miru, JA-C236-kiku, JA-C236-ongaku, JA-C236-ensou

\chapter{Films and Music}
\label{ch:ja-films-and-music}
\hlchaptermodality{236}

\begin{chapteropening}
\textbf{By the end of this chapter:} \emph{I can say what I watched and listened to.}

Complete the last lesson of chapter 236: I can say what I watched and listened to.
\end{chapteropening}

\section[\texorpdfstring{\ja{えいが} (eiga)}{えいが (eiga)}]{\ja{えいが} (eiga) — a film; a movie}

\label{lesson:JA-C236-eiga}



\begin{quote}
Before the new one: say the Japanese for judo, then the Japanese for a championship win.
\end{quote}

\subsection*{You'll want to know: \ja{えいが}}
\textbf{\ja{えいが}} — \emph{eiga} — \textquotedblleft{}a film; a movie\textquotedblright{}.

\begin{quote}
\textbf{\ja{せんしゅう} \ja{かぞくと} \ja{えいがに} \ja{いきました。}} — \emph{senshū kazoku to eiga ni ikimashita} — \textquotedblleft{}Last week I went to a film with my family.\textquotedblright{}
\end{quote}

\begin{cousinweb}
\textbf{\ja{えいが}} (\emph{eiga}) is Sino-Japanese (kango): two Chinese characters read with their on'yomi, readings borrowed from Chinese. They are \emph{ei}, \textquotedblleft{}reflect, project\textquotedblright{} (Middle Chinese \emph{'jaengH}) and \emph{ga}, \textquotedblleft{}picture\textquotedblright{} (\emph{hweaH}) (KANJIDIC2; Middle Chinese in Baxter's transcription).
\end{cousinweb}

\begin{grammarlens}[title={What you've built}]
One more word for everyday life at A2.
\end{grammarlens}

\subsection*{Your turn}
\begin{itemize}
\raggedright
  \item \emph{Say it:} \emph{eiga}
  \item \emph{Say it:} \emph{eiga}, once more
  \item \emph{Say it:} \emph{senshū kazoku to eiga ni ikimashita} — Last week I went to a film with my family.
  \item \emph{Recall:} say the Japanese for judo, then the Japanese for a championship win, then say \emph{eiga} again
\end{itemize}

\begin{tcolorbox}[breakable,colback=teal!4,colframe=teal!35!black,title={Before you move on}]
What does \ja{えいが} mean? (\textquotedblleft{}A film; a movie\textquotedblright{}.) Say it once more.
\end{tcolorbox}
\section[\texorpdfstring{\ja{みる} (miru)}{みる (miru)}]{\ja{みる} (miru) — to watch; to see}

\label{lesson:JA-C236-miru}



\begin{quote}
Before the new one: say the Japanese for a championship win, then the Japanese for a film.
\end{quote}

\subsection*{You'll want to know: \ja{みる}}
\textbf{\ja{みる}} — \emph{miru} — \textquotedblleft{}to watch; to see\textquotedblright{}.

\begin{quote}
\textbf{\ja{きのう} \ja{うちで} \ja{えいがを} \ja{みました。}} — \emph{kinō uchi de eiga o mimashita} — \textquotedblleft{}Yesterday I watched a film at home.\textquotedblright{}
\end{quote}

Polite \textbf{\ja{みます}}, \emph{mimasu}; an ichidan verb.

\begin{cousinweb}
\textbf{\ja{みる}} (\emph{miru}) is a native Japanese word (wago), inherited rather than borrowed from Chinese; its kanji was attached later as a kun'yomi (JMdict, KANJIDIC2).
\end{cousinweb}

\begin{grammarlens}[title={What you've built}]
One more word for everyday life at A2.
\end{grammarlens}

\subsection*{Your turn}
\begin{itemize}
\raggedright
  \item \emph{Say it:} \emph{miru}
  \item \emph{Say it:} \emph{miru}, once more
  \item \emph{Say it:} \emph{kinō uchi de eiga o mimashita} — Yesterday I watched a film at home.
  \item \emph{Recall:} say the Japanese for a championship win, then the Japanese for a film, then say \emph{miru} again
\end{itemize}

\begin{tcolorbox}[breakable,colback=teal!4,colframe=teal!35!black,title={Before you move on}]
What does \ja{みる} mean? (\textquotedblleft{}To watch; to see\textquotedblright{}.) Say it once more.
\end{tcolorbox}
\section[\texorpdfstring{\ja{きく} (kiku)}{きく (kiku)}]{\ja{きく} (kiku) — to listen; to hear; to ask}

\label{lesson:JA-C236-kiku}



\begin{quote}
Before the new one: say the Japanese for a film, then the Japanese for to watch.
\end{quote}

\subsection*{You'll want to know: \ja{きく}}
\textbf{\ja{きく}} — \emph{kiku} — \textquotedblleft{}to listen; to hear; to ask\textquotedblright{}.

\begin{quote}
\textbf{\ja{けさ} \ja{うちで} \ja{うたを} \ja{ききました。}} — \emph{kesa uchi de uta o kikimashita} — \textquotedblleft{}This morning I listened to songs at home.\textquotedblright{}
\end{quote}

Polite \textbf{\ja{ききます}}, \emph{kikimasu}; a godan verb.

\begin{cousinweb}
\textbf{\ja{きく}} (\emph{kiku}) is a native Japanese word (wago), inherited rather than borrowed from Chinese; its kanji was attached later as a kun'yomi (JMdict, KANJIDIC2).
\end{cousinweb}

\begin{grammarlens}[title={What you've built}]
One more word for everyday life at A2.
\end{grammarlens}

\subsection*{Your turn}
\begin{itemize}
\raggedright
  \item \emph{Say it:} \emph{kiku}
  \item \emph{Say it:} \emph{kiku}, once more
  \item \emph{Say it:} \emph{kesa uchi de uta o kikimashita} — This morning I listened to songs at home.
  \item \emph{Recall:} say the Japanese for a film, then the Japanese for to watch, then say \emph{kiku} again
\end{itemize}

\begin{tcolorbox}[breakable,colback=teal!4,colframe=teal!35!black,title={Before you move on}]
What does \ja{きく} mean? (\textquotedblleft{}To listen; to hear; to ask\textquotedblright{}.) Say it once more.
\end{tcolorbox}
\section[\texorpdfstring{\ja{おんがく} (ongaku)}{おんがく (ongaku)}]{\ja{おんがく} (ongaku) — music}

\label{lesson:JA-C236-ongaku}



\begin{quote}
Before the new one: say the Japanese for to watch, then the Japanese for to listen.
\end{quote}

\subsection*{You'll want to know: \ja{おんがく}}
\textbf{\ja{おんがく}} — \emph{ongaku} — \textquotedblleft{}music\textquotedblright{}.

\begin{quote}
\textbf{\ja{きのう} \ja{ともだちと} \ja{おんがくを} \ja{ききました。}} — \emph{kinō tomodachi to ongaku o kikimashita} — \textquotedblleft{}Yesterday I listened to music with a friend.\textquotedblright{}
\end{quote}

\begin{cousinweb}
\textbf{\ja{おんがく}} (\emph{ongaku}) is Sino-Japanese (kango): two Chinese characters read with their on'yomi, readings borrowed from Chinese. They are \emph{on}, \textquotedblleft{}sound\textquotedblright{} (Middle Chinese \emph{'im}) and \emph{gaku}, \textquotedblleft{}music\textquotedblright{} (\emph{ngaewk}) (KANJIDIC2; Middle Chinese in Baxter's transcription).
\end{cousinweb}

\begin{grammarlens}[title={What you've built}]
One more word for everyday life at A2.
\end{grammarlens}

\subsection*{Your turn}
\begin{itemize}
\raggedright
  \item \emph{Say it:} \emph{ongaku}
  \item \emph{Say it:} \emph{ongaku}, once more
  \item \emph{Say it:} \emph{kinō tomodachi to ongaku o kikimashita} — Yesterday I listened to music with a friend.
  \item \emph{Recall:} say the Japanese for to watch, then the Japanese for to listen, then say \emph{ongaku} again
\end{itemize}

\begin{tcolorbox}[breakable,colback=teal!4,colframe=teal!35!black,title={Before you move on}]
What does \ja{おんがく} mean? (\textquotedblleft{}Music\textquotedblright{}.) Say it once more.
\end{tcolorbox}
\section[\texorpdfstring{\ja{えんそう} (ensō)}{えんそう (ensō)}]{\ja{えんそう} (ensō) — a musical performance; a recital}

\label{lesson:JA-C236-ensou}



\begin{quote}
Before the new one: say the Japanese for to listen, then the Japanese for music.
\end{quote}

\subsection*{You'll want to know: \ja{えんそう}}
\textbf{\ja{えんそう}} — \emph{ensō} — \textquotedblleft{}a musical performance; a recital\textquotedblright{}.

\begin{quote}
\textbf{\ja{せんしゅう} \ja{いもうとの} \ja{えんそうを} \ja{ききました。}} — \emph{senshū imōto no ensō o kikimashita} — \textquotedblleft{}Last week I listened to my younger sister's recital.\textquotedblright{}
\end{quote}

\begin{cousinweb}
\textbf{\ja{えんそう}} (\emph{ensō}) is Sino-Japanese (kango): two Chinese characters read with their on'yomi, readings borrowed from Chinese. They are \emph{en}, \textquotedblleft{}perform\textquotedblright{} (Middle Chinese \emph{yenX}) and \emph{sō}, \textquotedblleft{}play music\textquotedblright{} (\emph{tsuwH}) (KANJIDIC2; Middle Chinese in Baxter's transcription).
\end{cousinweb}

\begin{grammarlens}[title={What you've built}]
One more word for everyday life at A2.
\end{grammarlens}

\subsection*{Your turn}
\begin{itemize}
\raggedright
  \item \emph{Say it:} \emph{ensō}
  \item \emph{Say it:} \emph{ensō}, once more
  \item \emph{Say it:} \emph{senshū imōto no ensō o kikimashita} — Last week I listened to my younger sister's recital.
  \item \emph{Recall:} say the Japanese for to listen, then the Japanese for music, then say \emph{ensō} again
\end{itemize}

\begin{tcolorbox}[breakable,colback=teal!4,colframe=teal!35!black,title={Before you move on}]
What does \ja{えんそう} mean? (\textquotedblleft{}A musical performance; a recital\textquotedblright{}.) Say it once more.
\end{tcolorbox}
